\chapter{Addestrare una rete neurale}
\label{cap:addestramento}

\begin{inbreve}
Con la backpropagation si sa calcolare il gradiente di ogni peso; restano le scelte che servono per addestrare
davvero un MLP: inizializzazione, minimi multipli, on-line/batch/mini-batch, learning rate e momentum, criteri di
arresto, overfitting e regolarizzazione, numero di unità (con il Cascade Correlation), rappresentazione di input e
output, validazione specifica per le reti. Il capitolo si chiude con pregi e difetti delle reti neurali e con un
banco di prova per verificare un'implementazione.
\end{inbreve}

\section{L'addestramento come percorso}

\begin{idea}[Un'interpretazione utile]
La backpropagation è un \textbf{percorso nello spazio dei pesi}, sulla superficie della loss. Il percorso dipende
dai dati, dal modello, dal \textbf{punto di partenza} (pesi iniziali), dalla \textbf{velocità} di convergenza e dal
\textbf{punto finale} (regola di arresto). Tutte queste scelte sono un \textbf{controllo della ricerca} nello spazio
delle ipotesi.
\end{idea}

L'algoritmo di base resta la discesa del gradiente: pesi iniziali piccoli, $\Delta\w = -\partial E/\partial\w$
calcolato con la backpropagation (fase forward fino alle uscite, fase backward per i delta),
$\w_{\mathit{new}} = \w+\eta\,\Delta\w$, ripetuto. Il modello è spesso \textbf{sovra-parametrizzato} e
l'ottimizzazione è \textbf{non convessa} e potenzialmente instabile: da qui le questioni di questo capitolo. Gli
\termine{iperparametri} sono i valori da fissare per eseguire l'addestramento. Le tecniche segnate con
\textbf{(!)} sono di uso comune e vanno usate in pratica; le altre sono varianti da conoscere. Le ``regole
empiriche'' trovate in rete vanno evitate, a meno che non siano dimostrate in letteratura scientifica o se ne sia
verificato sperimentalmente il vantaggio.

\section{Inizializzazione e minimi multipli}

\begin{idea}[(!) Inizializzazione]
Inizializzare i pesi con \textbf{valori casuali vicini a zero}, ad esempio nell'intervallo $[-0{,}7;\,+0{,}7]$
per dati standardizzati.
\end{idea}

Da evitare:
\begin{itemize}
  \item pesi \textbf{tutti nulli}: i delta delle unità nascoste sarebbero nulli (dipendono dai pesi verso
        l'uscita, \eqref{eq:delta-nascosta}) e la rete non imparerebbe;
  \item pesi \textbf{grandi}: le sigmoidi partono \textbf{saturate}, con derivate quasi nulle e apprendimento
        lentissimo;
  \item pesi \textbf{tutti uguali}: le unità nascoste restano \textbf{simmetriche}, ricevono gli stessi
        aggiornamenti e calcolano tutte la stessa cosa. La casualità serve a \textbf{rompere la simmetria}.
\end{itemize}
Si può tenere conto del \termine{fan-in} (numero di input di un'unità), ad esempio con un intervallo
$\text{range}\cdot2/\text{fan-in}$, ma non se il fan-in è troppo grande e non per le unità di uscita
(altrimenti i delta partono vicini a zero). Molto usata è l'inizializzazione di \textbf{Glorot e Bengio}: bias a 0
e pesi da una distribuzione uniforme in $[-1/a,+1/a]$ con $a = \sqrt{\text{fan-in}}$ (in varianti successive $a$
dipende anche dal fan-out). Una rete molto piccola può dipendere fortemente dall'inizializzazione.

\subsection{Minimi multipli}

La loss \textbf{non è convessa}: ci sono molti minimi locali e il risultato dipende dai pesi iniziali.

\begin{idea}[(!) Più inizializzazioni]
Provare diverse configurazioni iniziali casuali (ad esempio 5--10 o più \ingl{trial}) e riportare la
\textbf{media} e la \textbf{varianza} dei risultati. Come scegliere il modello finale tra le prove è discusso nel
\rif{sec:validazione-reti}; si possono anche usare tutti i modelli con una risposta \textbf{di comitato}
(\rif{sec:ensemble}).
\end{idea}

Fermarsi in un minimo locale con errore alto non è un grosso problema: lo si vede dall'errore finale e si può
ricominciare; spesso un \textbf{``buon'' minimo locale è sufficiente}. In ogni caso il minimo globale non è
l'obiettivo:
\begin{enumerate}
  \item nel ML \textbf{non si cerca né il minimo locale né quello globale di $\Remp$}, ma il minimo di $R$, che
        non si può calcolare; spesso ci si ferma in un punto che non ha nemmeno gradiente nullo;
  \item la rete è uno spazio delle ipotesi di dimensione variabile (\rif{sec:vc-reti}): durante l'addestramento
        l'errore di training scende verso il minimo mentre la rete diventa troppo complessa, e bisogna fermarsi
        \textbf{prima} di questa condizione di \ingl{overtraining}, indipendentemente dalla questione dei minimi.
\end{enumerate}

\section{On-line, batch e mini-batch}
\label{sec:minibatch}

Con la notazione della backpropagation, per il peso $w_{tu}$:
\[
\text{batch: } -\frac{\partial E}{\partial w_{tu}} = -\sum_{p=1}^{l}\frac{\partial E_p}{\partial w_{tu}} = \sum_{p=1}^{l}\delta_{p,t}\,o_{p,u},
\qquad
\text{on-line: } -\frac{\partial E_p}{\partial w_{tu}} = \delta_{p,t}\,o_{p,u}.
\]
Nel batch si somma su tutta l'epoca e poi si aggiorna (\rif{sec:batch-online}); nell'on-line si aggiorna dopo ogni
esempio. Il batch dà una stima \textbf{più accurata} del gradiente. Nell'on-line il gradiente di un esempio è
un'approssimazione \textbf{rumorosa} del gradiente complessivo (da cui il nome \termine{stochastic gradient
descent}, SGD): la discesa procede a zig-zag, e questo \textbf{può aiutare a sfuggire ai minimi locali}.

\begin{idea}[(!) Shuffling]
Nella versione on-line i pattern vanno presentati in \textbf{ordine diverso} a ogni epoca, di solito
\textbf{mescolandoli casualmente}, per evitare una deriva sistematica della discesa.
\end{idea}

Nel \termine{mini-batch} si divide ogni epoca in parti: si sommano i gradienti di $mb$ esempi ($1<mb<l$) prima di
aggiornare i pesi, e si ripete finché tutti i dati dell'epoca sono stati usati (\figrif{fig:minibatch}). Lo SGD
con mini-batch usa il campionamento casuale (shuffling) per una stima non distorta del gradiente; poiché il costo
di un aggiornamento dipende da $mb$ e non da $l$, è un fattore chiave per scalare a dataset di milioni di esempi
(con dati enormi o in streaming si può fare anche a meno delle epoche).

\begin{figure}[ht]
\centering
\begin{tikzpicture}[x=1cm, y=1cm]
\foreach \m in {1,...,4} {
  \draw[fill=grigiochiaro] (\m*2.3-2.3,0) rectangle ++(2.3,0.7) node[midway, font=\small] {$mb_\m$};
  \draw[->, rosso] (\m*2.3,0.75) -- ++(0,0.45) node[above, font=\scriptsize] {aggiornamento};
}
\draw[decorate, decoration={brace, mirror, amplitude=5pt}] (0,-0.1) -- (9.2,-0.1) node[midway, below=5pt, font=\footnotesize] {un'epoca: tutti gli $l$ esempi};
\end{tikzpicture}
\caption{Un'epoca divisa in quattro mini-batch: i pesi si aggiornano dopo ogni mini-batch.}
\label{fig:minibatch}
\end{figure}

\begin{table}[ht]
\centering\small
\caption{On-line, mini-batch e batch.}
\begin{tabularx}{0.95\linewidth}{XXX}
\toprule
On-line & Mini-batch (SGD) & Batch \\
\midrule
non sfrutta il parallelismo tra esempi & compromesso; si adatta alla memoria della GPU & stima più accurata del gradiente \\
può essere instabile & & più lento (aspetta tutti i dati) \\
il rumore può avere un effetto regolarizzante & & \\
\bottomrule
\end{tabularx}
\end{table}

\section{Il learning rate}
\label{sec:learning-rate}

\begin{itemize}
  \item \textbf{Batch}: gradiente più accurato, si può usare un $\eta$ più alto.
  \item \textbf{On-line}: più veloce ma più instabile, serve un $\eta$ più piccolo.
  \item In generale un $\eta$ \textbf{alto} è veloce ma può essere instabile, uno \textbf{basso} è stabile ma può
        essere troppo lento; con LMS si cerca tipicamente in $[0{,}01;\,0{,}5]$.
\end{itemize}

\begin{idea}[(!) Guardare sempre la curva di apprendimento]
Il grafico dell'errore di training durante l'addestramento mostra la \textbf{qualità} della discesa e la
\textbf{velocità} di convergenza (il valore assoluto dell'errore dipende anche dalla capacità del modello). Si cerca
una curva che scende rapidamente \textbf{e in modo liscio} (\figrif{fig:eta-curve}); una curva irregolare può
dipendere da batch troppo piccoli o da $\eta$ troppo alto, e va evitata. Qui si guardano solo i fallimenti
dell'ottimizzazione sul training, non la generalizzazione.
\end{idea}

\begin{figure}[ht]
\centering
\begin{tikzpicture}
\begin{axis}[stile, width=8.6cm, height=5cm, domain=0:100, samples=200, xmin=0, xmax=104, ymin=0, ymax=1.25,
  xtick=\empty, ytick=\empty, xlabel={epoche}, ylabel={errore di training}, legend pos=outer north east]
\addplot[thick, viola] {min(1.2, 0.55*exp(x/38))}; \addlegendentry{molto alto}
\addplot[thick, verde] {0.95-0.0008*x}; \addlegendentry{molto basso}
\addplot[thick, blu] {0.2+0.75*exp(-x/45)}; \addlegendentry{basso}
\addplot[thick, arancio] {0.42+0.55*exp(-x/5)}; \addlegendentry{alto}
\addplot[very thick, rosso] {0.08+0.87*exp(-x/14)}; \addlegendentry{buono}
\end{axis}
\end{tikzpicture}
\caption{Effetto del learning rate sulla curva di training (andamenti qualitativi): con $\eta$ molto alto la
discesa diverge, molto basso resta quasi ferma, basso scende lentamente, alto scende in fretta ma si ferma presto;
quello buono scende rapidamente e in modo liscio.}
\label{fig:eta-curve}
\end{figure}

\begin{idea}[(!) Usare la media dei gradienti]
Conviene usare la \textbf{media} dei gradienti sull'epoca (somma divisa per $l$): rende l'approccio uniforme
rispetto al numero di dati (è il \ingl{least mean squares}) ed equivale a usare $\eta/l$. Confrontando on-line e
batch, l'on-line richiede un $\eta$ molto più piccolo (dello stesso fattore $1/l$) per essere comparabile; nel
mini-batch, se si divide per $mb$, va ritarato $\eta$, oppure si divide sempre per $l$.
\end{idea}

\subsection{Momentum (!)}

\begin{equation}\label{eq:momentum}
\boxed{\;\Delta\w_{\mathit{new}} = -\eta\,\frac{\partial E(\w)}{\partial\w} + \alpha\,\Delta\w_{\mathit{old}},\qquad \w_{\mathit{new}} = \w + \Delta\w_{\mathit{new}}\;}
\end{equation}
con $0<\alpha<1$ (ad esempio $0{,}5$--$0{,}9$); qui $\eta$ è incluso in $\Delta\w$ e dopo ogni passo si salva
$\Delta\w_{\mathit{old}} = \Delta\w_{\mathit{new}}$. È il metodo della \textbf{``palla pesante''}: a ogni passo si
aggiunge una frazione dello spostamento precedente, come se i pesi avessero un'inerzia.
\begin{itemize}
  \item \textbf{Più veloce nei plateau}: se il gradiente ha lo stesso segno in iterazioni consecutive, i passi si
        sommano e il passo effettivo cresce.
  \item \textbf{Smorza le oscillazioni}: nelle direzioni in cui il gradiente cambia segno i contributi si
        compensano (\figrif{fig:momentum}).
  \item Per l'effetto inerzia permette di usare $\eta$ più alti.
\end{itemize}

\begin{figure}[ht]
\centering
\begin{tikzpicture}[scale=1.0]
\foreach \c in {0.4,0.8,1.2,1.6,2.0} \draw[grigio] (0,0) ellipse ({2.7*\c} and {0.42*\c});
\fill (0,0) circle (1.5pt);
\draw[->, thick, inchiostro] (-4.6,0.55) -- (-4.2,-0.5) -- (-3.6,0.42) -- (-3.1,-0.36) -- (-2.5,0.3) -- (-2.0,-0.24) -- (-1.4,0.17) -- (-0.9,-0.1);
\draw[->, very thick, rosso] (-4.6,0.55) .. controls (-3.8,-0.2) and (-2.4,-0.05) .. (-0.15,0.0);
\node[font=\footnotesize] at (-2.8,-1.15) {senza momentum: zig-zag tra le pareti};
\node[font=\footnotesize, rosso] at (-2.0,0.75) {con momentum};
\end{tikzpicture}
\caption{Una valle stretta della superficie d'errore (Hessiana mal condizionata): il gradiente puro rimbalza tra
le pareti, il momentum procede lungo la valle verso il minimo.}
\label{fig:momentum}
\end{figure}

Il momentum nasce per il batch. Nell'on-line $\Delta\w_{\mathit{old}}$ è quello dell'\textbf{esempio
precedente}, e il momentum diventa una \textbf{media mobile dei gradienti passati} che liscia i gradienti stocastici
dentro l'epoca (analogamente nel mini-batch): il significato è diverso dal batch, dove si combinano gradienti
calcolati sugli stessi esempi in epoche diverse.

Nel \termine{momentum di Nesterov} si applica prima il momentum ($\w+\alpha\Delta\w_{\mathit{old}}$), si valuta il
gradiente in questo punto intermedio, e poi si applica il $\Delta\w$ complessivo (momentum più nuovo gradiente) ai
pesi originali. Migliora la velocità di convergenza nel batch (non nel caso stocastico).

\subsection{Learning rate variabile e adattivo}

Con il mini-batch il gradiente \textbf{non tende a zero} vicino al minimo (per il rumore di campionamento), quindi
un $\eta$ fisso va evitato. Ad esempio lo si decresce linearmente fino all'iterazione $\tau$,
\[
\eta_s = \Big(1-\frac{s}{\tau}\Big)\eta_0 + \frac{s}{\tau}\,\eta_\tau,\qquad s\le\tau,
\]
e poi si usa $\eta=\eta_\tau$ costante e piccolo (ad esempio $\eta_\tau\approx1\%$ di $\eta_0$, $\tau$ qualche
centinaio di passi); $\eta_0$ si sceglie con il compromesso tra instabilità e blocco (prove preliminari o grid
search). Gli algoritmi a \textbf{learning rate adattivo} lo regolano automaticamente, anche separatamente
\textbf{per ogni parametro}: \textbf{AdaGrad}, \textbf{RMSProp}, \textbf{Adam} (molto popolare e spesso robusto con i
valori di default, ma da usare con cautela); si combinano con il momentum. Nelle reti con molti strati $\eta$ può
dipendere anche dallo strato o dal fan-in.

\subsection{Oltre lo SGD di base}

Lo SGD è efficiente (lineare nel numero di parametri) ma dipende dal numero di dati e di epoche. Esistono molte
euristiche: \textbf{Quickprop} (salta al minimo di una parabola che approssima localmente la loss),
\textbf{R-prop} (usa solo il \textbf{segno} del gradiente, utile quando il gradiente svanisce negli strati
profondi). I \textbf{metodi del secondo ordine} usano l'Hessiana (derivate seconde) per tenere conto della
curvatura (metodo di Newton), ma l'Hessiana è costosissima; ne esistono versioni approssimate (Gauss-Newton,
Levenberg-Marquardt, metodi \ingl{Hessian-free} come il gradiente coniugato, metodi quasi-Newton come BFGS).

\begin{attenzione}[Punti di sella]
Un punto di sella ha gradiente (quasi) nullo ma non è un minimo. In alta dimensione i punti di sella sono
esponenzialmente più numerosi dei minimi locali. Il metodo di Newton cerca i punti a gradiente nullo e quindi è
attratto dai punti di sella: può essere uno dei motivi per cui i metodi del secondo ordine non hanno sostituito la
discesa del gradiente per le reti neurali, mentre la discesa del gradiente empiricamente riesce a sfuggire ai punti
di sella.
\end{attenzione}

Gli approcci di base (quelli segnati con (!)), applicati correttamente, di solito bastano: l'esperienza con
tecniche note conta più di metodi sofisticati che non si sanno gestire. La tecnica di ottimizzazione conta solo se
ci sono ostacoli all'addestramento: \textbf{la validazione conta più dell'ottimizzazione sul training set}, che non
è l'obiettivo. Il punto di partenza è lo \textbf{SGD con momentum} (e regolarizzazione).

\section{Criteri di arresto}
\label{sec:arresto}

\begin{itemize}
  \item Il criterio di base è l'errore (errore medio $<\epsilon$): il migliore se si conosce la tolleranza dei dati
        (conoscenza dell'esperto), il che spesso non accade. Varianti: tolleranza massima invece che media; per la
        classificazione, il numero di errori.
  \item Criteri ``interni'': pesi che non cambiano più, gradiente quasi nullo (norma $<\epsilon$), errore che non
        diminuisce più in modo significativo in un'epoca (ad esempio meno dello 0,1\%). Possono essere
        \textbf{prematuri} (ad esempio con $\eta$ troppo piccolo), quindi si osservano $k$ epoche
        (\termine{patience}).
  \item \textbf{(!)} Fermarsi comunque dopo un numero eccessivo di epoche, per sfuggire a convergenze troppo lente,
        ma \textbf{non} fermarsi a un numero arbitrario di epoche fissato a priori (\rif{sec:validazione-reti}).
  \item Non necessariamente fermarsi con un errore di training molto basso: è il tema dell'overfitting.
\end{itemize}

\section{Overfitting e regolarizzazione}
\label{sec:reg-reti}

Tipicamente \textbf{non} si vuole il minimizzatore globale di $\Remp(\w)$, perché è probabilmente una soluzione in
overfitting: è la differenza fondamentale rispetto all'ottimizzazione pura. Lo scopo è il \textbf{controllo della
complessità}, tramite regolarizzazione (esplicita con una penalità, o implicita con l'early stopping) più la model
selection per trovare il compromesso.

\subsection{Come nasce l'overfitting in una rete}
\begin{itemize}
  \item Si parte con \textbf{pesi piccoli e casuali}.
  \item La mappatura input-output è quasi \textbf{lineare} (le sigmoidi lavorano nella zona lineare): il numero
        effettivo di parametri liberi, e la VC-dimension, è quasi quello di un'unità lineare.
  \item Man mano che l'ottimizzazione procede, le unità nascoste tendono a \textbf{saturare}, il numero effettivo
        di parametri e le non linearità aumentano: la VC-dimension cresce.
\end{itemize}

\begin{figure}[ht]
\centering
\begin{tikzpicture}
\begin{axis}[stile, width=8.4cm, height=4.8cm, domain=0:100, samples=150, xmin=0, xmax=104, ymin=0, ymax=1.05,
  xtick=\empty, ytick=\empty, xlabel={epoche}, ylabel={errore}, legend pos=north east]
\addplot[thick, blu] {0.05+0.9*exp(-x/18)}; \addlegendentry{training}
\addplot[thick, rosso] {0.25+0.7*exp(-x/14)+0.0035*max(0,x-30)}; \addlegendentry{validazione}
\draw[densely dashed, grigio] (axis cs:31,0) -- (axis cs:31,0.45);
\node[font=\footnotesize, grigio, align=center] at (axis cs:31,0.55) {zona in cui\\fermarsi};
\node[font=\footnotesize, grigio] at (axis cs:85,0.6) {overtraining};
\end{axis}
\end{tikzpicture}
\caption{Comportamento tipico della backpropagation (qualitativo): l'errore di training decresce sempre, quello
di validazione raggiunge un minimo e poi risale.}
\label{fig:early-stopping}
\end{figure}

La curva della \figrif{fig:early-stopping} va letta insieme al bound della SLT. Gli strumenti sono tre:
\begin{enumerate}
  \item \termine{early stopping}: usare un \textbf{validation set} per decidere quando fermarsi, cioè quando
        l'errore di validazione inizia a salire. È un criterio vago: bisogna osservare più epoche prima di decidere
        (patience) ed eventualmente tornare indietro (\ingl{backtracking}). Poiché il numero effettivo di parametri
        cresce durante l'addestramento, \textbf{fermarsi prima equivale a limitare la complessità effettiva};
        \label{sec:early-stopping}
  \item \textbf{regolarizzazione} della loss, lo strumento principale;
  \item metodi di \textbf{pruning} (eliminazione di pesi o unità).
\end{enumerate}

\subsection{Regolarizzazione (!)}

Si usa la penalità di Tikhonov (\rif{sec:tikhonov}) su \textbf{tutti i pesi della rete}:
\[
\mathrm{Loss}(\w) = \sum_p\big(d_p - o(\x_p)\big)^2 + \lambda\norm{\w}^2,\qquad \norm{\w}^2 = \sum_iw_i^2 ,
\]
con l'effetto di \textbf{weight decay} \eqref{eq:weight-decay}: si aggiunge circa $2\lambda w$ al gradiente di ogni
peso (il 2 si può omettere). $\lambda$ è in genere piccolo (ad esempio $0{,}01$) e si sceglie con la model
selection. Applicata a un modello lineare è la ridge regression; esistono penalità più sofisticate (ad esempio per
la \ingl{weight elimination}). Nei risultati \textbf{(!)} si riporta l'errore (il termine sui dati), non la loss.

\begin{figure}[ht]
\centering
\begin{tikzpicture}
\begin{groupplot}[group style={group size=2 by 1, horizontal sep=0.45cm}, scatola, width=6.4cm, height=5.8cm,
  xmin=-3, xmax=3.8, ymin=-2.6, ymax=3.4, xtick=\empty, ytick=\empty]
\nextgroupplot[title={\small senza weight decay}]
\addplot[only marks, mark=o, mark size=1.4pt, rosso!80!black, thin] table[x=x, y=y, restrict expr to domain={\thisrow{c}}{0:0}] {dati/htf_punti.dat};
\addplot[only marks, mark=*, mark size=1.2pt, blu] table[x=x, y=y, restrict expr to domain={\thisrow{c}}{1:1}] {dati/htf_punti.dat};
\addplot[thick, viola, densely dashed] table {dati/htf_bayes.dat};
\addplot[thick, inchiostro] table {dati/nn_senza.dat};
\nextgroupplot[title={\small con weight decay}]
\addplot[only marks, mark=o, mark size=1.4pt, rosso!80!black, thin] table[x=x, y=y, restrict expr to domain={\thisrow{c}}{0:0}] {dati/htf_punti.dat};
\addplot[only marks, mark=*, mark size=1.2pt, blu] table[x=x, y=y, restrict expr to domain={\thisrow{c}}{1:1}] {dati/htf_punti.dat};
\addplot[thick, viola, densely dashed] table {dati/htf_bayes.dat};
\addplot[thick, inchiostro] table {dati/nn_wd.dat};
\end{groupplot}
\end{tikzpicture}
\caption{Una rete con 10 unità nascoste sul problema a due classi, senza e con weight decay (nero; tratteggiato il
confine di Bayes). Con la regolarizzazione il confine è più regolare e l'errore di test scende da
$\valnnsenzaTest$ a $\valnnwdTest$ (errore di Bayes $\valbayesErr$).}
\label{fig:weight-decay}
\end{figure}

\begin{attenzione}[La regolarizzazione non serve a stabilizzare la convergenza]
La regolarizzazione serve a \textbf{controllare la complessità} del modello (la VC-dimension, legata al numero e
al valore dei pesi), non a rendere più stabile la curva di apprendimento. Se la curva è instabile il problema è
altrove (learning rate, dimensione dei batch, \dots).
\end{attenzione}

\begin{idea}[Early stopping o regolarizzazione?]
L'early stopping è un approccio \textbf{empirico}, che richiede un validation set e sacrifica una parte dei dati.
La regolarizzazione nella loss è un approccio \textbf{fondato}: permette alla curva di validazione di seguire
quella di training (purché $\lambda$ non sia così alto da dare underfitting), e in quel caso non serve l'early
stopping e ci si può fermare a convergenza. Si possono usare entrambi: l'early stopping non interviene se l'errore
di validazione non sale.
\end{idea}

Alcuni dettagli:
\begin{itemize}
  \item spesso il \textbf{bias $w_0$ è escluso} dal regolarizzatore (includerlo renderebbe i risultati dipendenti
        da traslazioni o scalature del target), oppure ha un proprio coefficiente;
  \item la regolarizzazione si applica tipicamente nella versione batch; nell'on-line e nel mini-batch la penalità
        viene applicata a ogni passo, cioè molte volte per epoca, e per un confronto equo con il batch si userebbe
        $\lambda\cdot mb/l$ (se $\lambda$ si sceglie con la model selection, questo avviene automaticamente);
  \item un'altra tecnica di regolarizzazione è il \textbf{dropout} (\rif{sec:dropout}).
\end{itemize}

\subsection{Momentum e weight decay insieme}

Con la loss regolarizzata, separando i due termini e usando $\eta$ solo per il termine d'errore (come nel
\rif{sec:tikhonov}), la regola con momentum diventa
\[
\Delta w_{tu} = -\eta\,\frac{\partial\,\mathrm{Loss}}{\partial w_{tu}} + \alpha\,\Delta w_{tu}^{\mathit{old}}
\;\doteq\; \eta\,\delta_t\,o_u - \lambda\,w_{tu} + \alpha\,\Delta w_{tu}^{\mathit{old}},\qquad
w_{tu}^{\mathit{new}} = w_{tu} + \Delta w_{tu},
\]
così $\eta$ non moltiplica $\lambda$ e cambiando l'uno non si cambia l'altro. Spesso il momentum include anche la
parte di $\lambda$; per rendere \textbf{indipendenti tutti e tre gli iperparametri} $\eta$, $\alpha$, $\lambda$ si
può separare il weight decay dalla memoria del momentum:
\[
\Delta w_{tu} = \eta\,\delta_t\,o_u + \alpha\,\Delta w_{tu}^{\mathit{old}},\qquad
w_{tu}^{\mathit{new}} = w_{tu} + \Delta w_{tu} - \lambda\,w_{tu}.
\]
Si divide per $l$ solo il gradiente di $E$, dove c'è la somma sugli esempi.

\section{Numero di unità nascoste}
\label{sec:unita-nascoste}

Il numero di unità è legato al controllo della complessità, alla dimensione dell'input e a quella del training set:
in generale è un problema di \textbf{model selection}, e il numero di unità si sceglie insieme a $\lambda$ con la
cross-validation. Troppe poche unità danno underfitting, troppe possono dare overfitting; il numero può essere più
alto se si usa un'adeguata regolarizzazione. Alternative alla ricerca manuale sono gli \textbf{approcci
costruttivi}, che partono da una rete piccola e aggiungono unità, e i \textbf{metodi di pruning}, che partono da
reti grandi ed eliminano pesi o unità.

\subsection{Cascade Correlation}

Il \termine{Cascade Correlation} (CC) è l'approccio costruttivo più noto, per regressione e classificazione.
Apprende \textbf{sia i pesi sia la topologia}: lo spazio delle ipotesi ha dimensione decisa dall'algoritmo, e a
ogni passo si addestra una sola unità.

\begin{algoritmo}[Cascade Correlation]
\begin{enumerate}
  \item Si parte da una rete $N_0$ \textbf{senza unità nascoste} e la si addestra. Se non risolve il problema si
        passa a $N_1$.
  \item In $N_1$ si aggiunge un'unità nascosta i cui pesi in ingresso sono addestrati per \textbf{massimizzare la
        correlazione} tra la sua uscita e l'\textbf{errore residuo} della rete precedente.
  \item Dopo l'addestramento i pesi in ingresso della nuova unità vengono \textbf{congelati} e si riaddestrano i
        pesi rimanenti (lo strato di uscita).
  \item Se la rete non risolve il problema si aggiungono altre unità, ciascuna collegata a \textbf{tutti gli input e
        a tutte le unità nascoste precedenti} (da qui ``a cascata'').
  \item Si continua finché l'errore residuo soddisfa il criterio di arresto.
\end{enumerate}
\end{algoritmo}

\begin{figure}[ht]
\centering
\begin{tikzpicture}[x=1cm, y=1cm]
\foreach \n/\xs in {1/0, 2/4.6, 3/9.2} {
  \begin{scope}[xshift=\xs cm]
    \node[neur] (a) at (0,0) {$x_1$}; \node[neur] (b) at (0,-1) {$x_2$};
    \node[neurp] (o) at (3,-0.5) {$o$};
    \draw[->] (a) -- (o); \draw[->] (b) -- (o);
    \node[neurp] (h1) at (1.2,0.9) {$h_1$};
    \draw[->] (a) -- (h1) node[midway, left, font=\scriptsize] {$*$}; \draw[->] (b) -- (h1);
    \draw[->] (h1) -- (o);
    \ifnum\n>1
      \node[neurp] (h2) at (2.0,1.6) {$h_2$};
      \draw[->] (a) -- (h2); \draw[->] (b) to[bend left=10] (h2); \draw[->] (h1) -- (h2) node[midway, above, font=\scriptsize] {$*$};
      \draw[->] (h2) -- (o);
    \fi
    \ifnum\n>2
      \node[neurp] (h3) at (2.8,2.3) {$h_3$};
      \draw[->] (a) to[bend left=20] (h3); \draw[->] (h1) -- (h3); \draw[->] (h2) -- (h3) node[midway, above, font=\scriptsize] {$*$};
      \draw[->] (h3) -- (o);
    \fi
    \node[font=\footnotesize] at (1.5,-1.6) {\n\ unità nascost\ifnum\n=1 a\else e\fi};
  \end{scope}
}
\end{tikzpicture}
\caption{Crescita di una rete Cascade Correlation: ogni nuova unità riceve gli input e le uscite di tutte le unità
nascoste precedenti; i pesi segnati con $*$ (in ingresso alle unità già inserite) restano congelati.}
\label{fig:cascade}
\end{figure}

L'algoritmo alterna la minimizzazione dell'errore totale (LMS, ad esempio backpropagation sullo strato di uscita) e
la \textbf{massimizzazione della covarianza} $S$ tra l'uscita $o_p$ della candidata e l'errore residuo
$E_{p,k} = o_{p,k}-d_{p,k}$ di ogni uscita $k$:
\[
S = \sum_k\Big\lvert\sum_p\big(o_p-\bar o\big)\big(E_{p,k}-\bar E_k\big)\Big\rvert .
\]
Il gradiente, usando $d\abs{f}/dx = \sign(f)\,df/dx$ e trascurando la dipendenza delle medie da $o_p$, è
\[
\frac{\partial S}{\partial w_j} = \sum_k\sign(S_k)\sum_p\big(E_{p,k}-\bar E_k\big)\,f'(\net_{p,h})\,I_{p,j},
\]
dove $h$ è la candidata, $S_k$ il termine $k$-esimo della somma e $I_{p,j}$ il suo input $j$-esimo; si fa
\textbf{ascesa} del gradiente, perché si massimizza $S$.
\begin{itemize}
  \item Il ruolo delle unità nascoste è \textbf{ridurre l'errore residuo}: ogni unità risolve un sotto-problema e
        diventa un \textbf{rilevatore di feature permanente}.
  \item \textbf{Pool di candidate}: la superficie di $S$ ha molti massimi, quindi si addestra un gruppo di
        candidate e si sceglie la migliore.
  \item È un algoritmo \textbf{greedy}: converge facilmente e può trovare un numero minimo di unità, ma può andare
        in overfitting (serve regolarizzazione).
  \item La loss non deve essere necessariamente ``max $S$'' (per la regressione si può usare LMS).
\end{itemize}

\section{Rappresentazione di input e output}
\label{sec:input-output}

\textbf{Input.} Il pre-processing ha un grande effetto.
\begin{itemize}
  \item \textbf{Standardizzazione} (spesso utile): per ogni feature media 0 e deviazione standard 1,
        $(v-\text{media})/\text{dev.\,std}$.
  \item \textbf{Rescaling} in $[0,1]$: $(v-\min)/(\max-\min)$.
  \item Input categorici: codifica \textbf{1-of-$K$} (\rif{sec:dati}).
  \item Dati mancanti: attenzione, \textbf{0 non significa ``nessun input''} se 0 è nel range dei valori.
\end{itemize}

\textbf{Output.}
\begin{itemize}
  \item \textbf{(!) Regressione}: unità di uscita \textbf{lineari}.
  \item \textbf{Classificazione}: una unità (binaria) o codifica 1-of-$K$ con più uscite. Con la sigmoide si sceglie
        la soglia per assegnare la classe (eventualmente con zona di rifiuto); con 1-of-$K$ vince l'uscita più alta.
        Spesso la \textbf{tanh} (simmetrica) apprende più in fretta. Come target si può usare $0{,}9$ invece di 1
        (e $-0{,}9$ o $0{,}1$ invece di $-1$ o 0) per evitare la convergenza asintotica verso la saturazione; il
        range dei target deve comunque stare nel range delle uscite.
  \item \termine{Softmax} per target 0/1 multi-classe: le uscite sommano a 1 e si interpretano come probabilità
        $p(\text{classe}=k\mid\x)$:
        \[
        o_k(\x) = \frac{e^{\net_k}}{\sum_{j=1}^{K}e^{\net_j}} .
        \]
  \item Come loss alternativa si può usare la \termine{cross-entropy} (stima di massima verosimiglianza); per
        un'unità:
        \[
        -\sum_{i\in TR}\Big\{d_i\log\big(\out(\x_i)\big) + (1-d_i)\log\big(1-\out(\x_i)\big)\Big\}.
        \]
\end{itemize}

\section{Validazione di una rete neurale}
\label{sec:validazione-reti}

Gli schemi generali sono nel capitolo~\ref{cap:validazione}; per le reti ci sono tre questioni specifiche.

\subsection{Non fissare il numero di epoche}
\begin{attenzione}[Errore frequente]
Evitare di fermare l'addestramento dopo un \textbf{numero fisso e arbitrario di epoche}.
\end{attenzione}
Se è piccolo ci si ferma \textbf{troppo presto} (underfitting, o si favoriscono i learning rate alti che richiedono
meno epoche); se è grande \textbf{troppo tardi} (overfitting o tempo perso); lo stesso valore non può andare bene
per tutte le configurazioni e le iterazioni di una K-fold CV; e non si saprebbe quale numero usare nel
riaddestramento finale. Anche selezionare il numero di epoche con la model selection è impreciso. Con un sistema
\textbf{ben regolarizzato} funzionano bene i criteri del \rif{sec:arresto} (soglia sulla diminuzione dell'errore o
sul gradiente), oppure l'early stopping se si può usare un validation set.

\subsection{Early stopping e cross-validation}
Il criterio di early stopping fa \textbf{parte della model selection} (non va applicato sul test set) e in linea di
principio richiede un VL ogni volta che si addestra il modello. Il problema è scegliere il punto di arresto del
\textbf{riaddestramento finale}, se ogni fold della CV si è fermato a un'epoca diversa:
\begin{itemize}
  \item prendere la \textbf{media} delle epoche sui fold è rischioso: con più dati nel riaddestramento il punto di
        arresto ottimo cambia (ad esempio, in on-line, più dati per epoca significano più aggiornamenti per epoca e
        servono meno epoche);
  \item \textbf{molto meglio}: prendere la media dell'\textbf{errore di training} nel punto migliore di validazione e
        usarla come soglia nel riaddestramento, per raggiungere lo stesso livello di fitting;
  \item altrimenti serve un VL ogni volta: CV con early stopping attivo in ogni fold e, trovati i migliori
        iperparametri, riaddestramento sul dataset ricomposto usando ancora una parte dei dati come VL.
\end{itemize}
In alternativa all'early stopping si usano \textbf{modelli regolarizzati} con criteri di arresto standard; la
regolarizzazione riduce comunque la dipendenza dall'early stopping.

\subsection{Inizializzazioni diverse e modello finale}
Inizializzazioni diverse portano a \textbf{modelli diversi}. Come scegliere il modello finale (scelta da
dichiarare quando si riportano i risultati)?
\begin{itemize}
  \item Si calcolano \textbf{media e varianza} dell'errore sulle prove: un'esecuzione qualsiasi cadrà in
        quell'intervallo, quindi il rischio è sotto controllo, a meno di una varianza alta (che è comunque un
        segnale negativo, \rif{sec:decomposizione}).
  \item Si usano \textbf{tutti} i modelli con un \textbf{ensemble} (media delle uscite o voto).
  \item Se se ne sceglie \textbf{uno}, è una \textbf{model selection} e serve un VL: il migliore (errore minimo sul
        VL), la mediana (scelta più cauta), oppure uno casuale (che non richiede un VL).
  \item Si possono scartare i casi ``sfortunati'' guardando solo l'errore di training.
\end{itemize}

\section{Pregi e difetti delle reti neurali}

\begin{idea}[Pregi]
\begin{itemize}
  \item Metodo flessibile, che estende i metodi lineari alla stima di funzioni arbitrarie e costruisce
        un'\textbf{ipotesi esplicita} (a differenza del K-NN).
  \item Spazio delle ipotesi continuo e ricco, loss differenziabile: addestramento con la discesa del gradiente.
  \item Approssimatori universali; gestiscono rumore e dati incompleti con degrado graduale; dati continui e
        discreti, regressione e classificazione.
  \item Un paradigma ampio (anche apprendimento non supervisionato, memorie associative, approcci stocastici).
  \item \textbf{Rappresentazione distribuita} della conoscenza, compatta in matrici di pesi; tolleranza ai guasti;
        adattamento ad ambienti non stazionari con l'apprendimento on-line (\ingl{continual learning}, con il
        dilemma stabilità-plasticità).
\end{itemize}
\end{idea}

\begin{attenzione}[Punti critici]
\begin{itemize}
  \item \textbf{Black box}: è difficile interpretare la conoscenza acquisita (tema della \ingl{explainable AI}).
        Approcci: proiezione o clustering delle rappresentazioni interne, analisi di sensitività,
        \ingl{Layer-wise Relevance Propagation} (mappe di calore sulle parti dell'input rilevanti per la
        decisione), estrazione di regole simboliche.
  \item \textbf{Progettazione dell'architettura}: servono metodi costruttivi o di pruning, o approcci diversi (SVM).
  \item \textbf{Dipendenza dal processo di addestramento}: varianti e iperparametri influenzano la soluzione.
  \item Non ideali con \textbf{dati mancanti} (si usano rimozione, imputazione, \dots).
\end{itemize}
\end{attenzione}

Le reti neurali sono indicate quando l'input è ad alta dimensione (discreto o reale), il task è di classificazione
o regressione, i dati possono essere rumorosi, il tempo di addestramento non è critico, la forma della funzione
target è sconosciuta, la leggibilità del risultato non è critica e il calcolo dell'uscita deve essere veloce.

\section{Verificare un'implementazione: il benchmark MONK}

Verificare la correttezza di un'implementazione è difficile: una rete spesso ``funziona'' (anche se male)
nonostante piccoli errori. Un primo collaudo è il dataset \textbf{MONK}:
\begin{itemize}
  \item 3 task di classificazione binaria piccoli e artificiali: una piccola rete con poche unità raggiunge
        accuratezze altissime (fino al 100\%) in poco tempo;
  \item input: 6 variabili con 2, 3 o 4 valori simbolici ciascuna, quindi codifica \textbf{1-of-$k$} e
        \textbf{17 unità di input};
  \item il test set include il training set (cattiva pratica, qui irrilevante); MONK3 contiene etichette
        rumorose e può andare in overfitting.
\end{itemize}

\begin{idea}[Consigli per le prime prove]
\begin{itemize}
  \item Un'unità di uscita, con target e uscite coerenti ($0/1$ o $-1/+1$), ad esempio con tanh.
  \item \textbf{Codifica one-hot degli input}: dimenticarla è un errore frequentissimo.
  \item Batch con gradiente diviso per il numero di esempi: converge in qualche centinaio di epoche.
  \item Poche unità nascoste (2--5); provare intervalli di inizializzazione piccoli (o basati sul fan-in), perché
        con poche unità il risultato è sensibile all'inizializzazione.
  \item Con approcci on-line usare lo \textbf{shuffling}: gli esempi sono ordinati per classe.
\end{itemize}
Buoni risultati sul MONK \textbf{non garantiscono} la correttezza; risultati cattivi indicano \textbf{sicuramente}
che codice o impostazioni vanno rivisti. Nelle curve conta la \textbf{forma} (lisce, convergenti), non i valori.
Per verificare la backpropagation si può usare un esempio svolto passo per passo (attenzione se omette i bias) o
una libreria affidabile come confronto.
\end{idea}

\begin{puntichiave}
\begin{itemize}
  \item Pesi iniziali piccoli e casuali; più inizializzazioni, riportando media e varianza.
  \item On-line (rumoroso, $\eta$ piccolo, shuffling), batch (preciso), mini-batch (compromesso); gradiente medio.
  \item Momentum $\Delta\w = -\eta\nabla E + \alpha\Delta\w_{\mathit{old}}$: accelera nei plateau, smorza le
        oscillazioni.
  \item Overfitting: la VC-dimension effettiva cresce durante il training; rimedi: early stopping,
        regolarizzazione (weight decay), pruning.
  \item Uscite lineari per la regressione; sigmoide/tanh o softmax per la classificazione.
  \item Mai un numero fisso di epoche; l'early stopping è parte della model selection.
\end{itemize}
\end{puntichiave}

\begin{domande}
\begin{enumerate}
  \item Come si inizializzano i pesi e perché?
  \item Differenze tra on-line, batch e mini-batch.
  \item Ruolo del learning rate; cos'è il momentum e perché funziona?
  \item Serve trovare il minimo globale? Perché?
  \item Come si manifesta l'overfitting in una rete neurale? Cos'è l'early stopping?
  \item Scrivere la loss con weight decay e la regola di aggiornamento con momentum. La regolarizzazione migliora la
        stabilità della convergenza?
  \item Descrivere il Cascade Correlation.
  \item Come si rappresentano input e output per regressione e classificazione?
  \item Perché non bisogna fermare l'addestramento dopo un numero fisso di epoche? Come si gestisce l'early stopping
        con la cross-validation?
\end{enumerate}
\end{domande}
